\documentclass[journal=jprobs,manuscript=suppinfo,layout=traditional]{achemso}
\usepackage[T1]{fontenc}
\usepackage[utf8]{inputenc}
\usepackage{graphicx}
\usepackage{amsmath}
\usepackage{booktabs}
\usepackage{lmodern}
\SectionNumbersOn
\author{Murat G\"ok}
\affiliation{Department of Computer Engineering, Faculty of Engineering, Yalova University, 77200 Yalova, T\"urkiye}
\email{murat.gok@yalova.edu.tr}
\title{Supporting Information: Condition identity and label weighting govern structure--response inference from reused site-resolved PTM data}
\sloppy
\begin{document}
\maketitle

\section{Protein-level comparison with bulk responses}

The source study reports bulk CML-ELISA responses for casein, haemoglobin and
collagen at three modification levels each. Figure~\ref{fig:s1} summarises the
comparison discussed in the main text.

The observed ordering is stable: casein exceeds haemoglobin, which exceeds
collagen, at all three levels (Figure~\ref{fig:s1}A), so the comparison does not
fail on the measurement side. One value, casein at the low modification level,
exceeds unity (1.34, SD 0.61) and is therefore not physically interpretable as a
repair fraction; we retain it as reported and note that it does not affect the
ordering.

Across the 96 of 192 pipeline configurations in which all three scaffolds were
represented, seven recovered the observed ordering (Figure~\ref{fig:s1}B),
against a chance expectation of one in six. The two most frequent predicted
orderings place haemoglobin or collagen first. These configurations are nested
variants of a single pipeline and are strongly correlated, so the fraction is
descriptive and not a significance test.

The prespecified configuration could not be evaluated at all because collagen
contributed no lysine meeting both the structure-confidence and
applicability-domain criteria (Figure~\ref{fig:s1}C). Of 69 collagen lysines,
none met the strict structure-quality criterion and one fell inside the
applicability domain; for casein the corresponding counts were 8 and 18 of 58,
and for haemoglobin 24 and 12 of 24.

With three scaffolds, an exact-ordering test cannot yield a $P$ value below one
in six, so this comparison could not have supported a transfer claim under any
outcome. It is reported as an applicability stress test.

\begin{figure}[htbp]
\centering
\includegraphics[width=\textwidth]{figureS1_cross_scaffold.pdf}
\caption{Protein-level comparison. (A) Observed bulk repair fraction by scaffold
and modification level, with standard deviations over four replicates; the
dotted line marks complete repair. (B) Distribution of predicted scaffold
orderings across the 96 complete pipeline configurations; the observed ordering
is highlighted and the dashed line marks the chance expectation. Cas, casein;
Hb, haemoglobin; Col, collagen. (C) Lysines per scaffold, those meeting the
strict structure-quality criterion, and those inside the applicability domain.}
\label{fig:s1}
\end{figure}

\section{Response dispersion by substrate load}

\begin{table}[htbp]
\caption{Response distribution for the site sets analysed in the main text.}
\label{tbl:s1}
\begin{tabular}{lrrrrrrr}
\toprule
Load (ng) & $n$ & Mean & SD & Min & Max & IQR & $\le 0.10$ / $\ge 0.90$ \\
\midrule
250  & 22 & 0.547 & 0.357 & 0.00 & 0.96 & 0.574 & 5 / 5 \\
500  & 33 & 0.602 & 0.310 & 0.00 & 1.00 & 0.509 & 3 / 7 \\
1000 & 37 & 0.479 & 0.304 & 0.00 & 0.89 & 0.533 & 7 / 0 \\
2000 & 38 & 0.455 & 0.298 & 0.00 & 0.97 & 0.492 & 6 / 1 \\
\bottomrule
\end{tabular}
\end{table}

Dispersion is comparable across loads, with standard deviations spanning
0.298--0.357 and interquartile ranges 0.492--0.574, a difference of at most
twenty percent between the widest and narrowest condition. Mean response
declines with increasing substrate load, the direction expected when enzyme
concentration is held constant.

\section{Descriptor extraction and validation}

Structural descriptors were extracted for all 44 lysines appearing in the union
across loads. The descriptor table distributed with the original analysis covers
only the 33 sites of the 500~ng condition. To verify that our extraction
reproduces that pipeline, descriptors were recomputed for those 33 sites and
compared field by field. All thirteen fields agreed to floating-point precision,
the largest absolute difference being $1.4\times10^{-14}$ (residue
solvent-accessible surface area). The eleven additional lysines (positions 36,
44, 151, 160, 228, 235, 248, 285, 340, 544 and 548) are ordinary lysines in
chain A of PDB 4F5S and yield descriptors under the same rules.

\end{document}
